\documentclass[10pt,a4paper]{amsart}
\usepackage{amsmath,amssymb,amsthm,graphicx}
\usepackage{amsmath}
\usepackage{amsfonts}
\usepackage{amssymb}
\usepackage{graphicx}
\usepackage{float}
\usepackage{caption}
\usepackage{pgf,tikz}
\usetikzlibrary{arrows}
\usepackage{verbatim}
\usepackage{hyperref}

\definecolor{ffttww}{rgb}{1,0.2,0.4}
\definecolor{uuuuuu}{rgb}{0.27,0.27,0.27}
\definecolor{cqcqcq}{rgb}{0.75,0.75,0.75}
\definecolor{wwqqcc}{rgb}{0.4,0,0.8}

\newcommand{\R}{\mathbb{R}}
\newcommand{\N}{\mathbb{N}}
\newcommand{\HH}{\mathcal{H}}
\newcommand{\RR}{\mathcal{R}}
\newcommand{\A}{\mathcal{A}}
\newcommand{\B}{\mathcal{B}}
\newcommand{\est}{\emptyset}
\newtheorem{theorem}{Theorem}
\newtheorem{proposition}{Proposition}
\newtheorem{corollary}{Corollary}
\newtheorem{lemma}{Lemma}
\newtheorem{definition}{Definition}
\newtheorem{jel}{Notation}
\newtheorem{conj}{Conjecture}
\title{Exact polychromatic parameter of unions of s axis-parallel strips in the plane}
\author{Balázs Bursics, Bence Csonka, Luca Szepessy}
\date{arxiv:2307.12154v3; CC BY 4.0}
\begin{document}
\maketitle
\noindent\textit{Editorial note.} The counted statement is printed in the introduction of the article, while its proof stands in the Proofs section; the proof environment carries the author's own attribution, Proof of Theorem s\_strip, which is the sole pairing used here. No other proof of this theorem exists in the article.\par\medskip
\noindent\textit{Editorial note.} The article's two transported PNG figures (abra1\_en.png, abra3.png) belong to the separate strip construction and are not part of this item, so no asset is declared and no figure is deleted from the declared spans.\par\medskip
\noindent\textit{Editorial note.} The retained preamble is the native macro and theorem-environment block of the article; font, spacing, caption and penguin-drawing packages of the original file are dropped, and the statement, proof and context spans need none of them.\par\medskip
\noindent\textit{Editorial note.} Only the counted theorem is a counting unit. The printed corollary that follows it at source lines 257-259, which specialises the identity to two and three colours, is a direct corollary and is not counted; Theorem strip, Theorem kereszt, Theorem vegtelen\_sorozat and Theorem veges\_sorozat are uncounted neighbours as well.\par\medskip
\noindent\textit{Editorial scope.} One counted exact parameter identity for unions of s axis-parallel strips, m\_\{A\_s\}(k) = s*m\_A(k)-s+1, with both directions of the article's complete proof. The printed corollary specialising it, the primal strip and bottomless non-existence theorems, the A-plus and A-star bounds and both arithmetic-progression theorems are not counted.\par\medskip
\noindent\textit{Allowed standard prerequisites (editorial).}\begin{itemize}
\item Polychromatic colourings and the parameter m\_H(k). A polychromatic k-colouring of a hypergraph puts a vertex of every one of the k colour classes into every hyperedge; m\_H(k) is the least m such that every member of the family has a polychromatic k-colouring after all edges of size below m are deleted, and it is infinity when no such m exists.. source lines 120 and 123-129, retained as context. Supplies both sides of the counted identity: the quantity m\_A(k) that is being related to m\_\{A\_s\}(k), and the meaning of an edge that is not polychromatic.
\item Definitional unpacking of the extremal value m\_A(k). For every m below m\_A(k) there is a member of A together with a colouring under which no k-colouring of its m-edges is polychromatic; consequently, just below the exact value, such a member has a non-polychromatic edge of size at least m\_A(k)-1, while the colouring witnessing m\_A(k) makes every edge of size m\_A(k) polychromatic.. source lines 680 and 682, read off the definition at source line 128. The counted proof takes one such extremal member and its non-polychromatic edge of size m\_A(k)-1 and reuses the colouring of the reverse direction.
\item Translated copies of a planar strip hypergraph along a diagonal. Finitely many translated copies of a finite planar point set can be placed along a diagonal so that every subset captured by an axis-parallel strip in the original set is captured by an axis-parallel strip in the copy, and so that any s of the copies lie inside the union of s axis-parallel strips.. source line 680. Realises the k(s-1)+1 copies used in the lower-bound direction and lets s extracted bad edges be read as one edge of the family A\_s.
\item Elementary finite counting and the pigeonhole principle. Among k(s-1)+1 objects coloured with k colours some colour occurs at least s times; and s(m-1)+1 points distributed over s sets leave one set with at least m points.. source lines 680 and 682. Produces the s copies with a common missing colour in the lower-bound direction and the single strip with m\_A(k) points in the upper-bound direction.
\end{itemize}\medskip
A {\it polychromatic k-coloring} of a hypergraph $H$ is a $k$-coloring of its vertex set such that every hyperedge contains a vertex of all $k$ color classes. Observe that a polychromatic 2-coloring is the same as the usual proper 2-coloring of hypergraphs, where we require that no edge is monochromatic. By merging $j$ color classes of a polychromatic $k$-coloring we get a polychromatic $(k-j+1)$-coloring, so the condition of monochromatic $k$-colorability becomes stricter as $k$ increases. A trivial necessary condition for the existence of a polychromatic $k$-coloring is that all edges of $H$ must be of size at least $k$.

\begin{definition} Denote by $H_{\ge m}$ or simply $H_m$ the hypergraph obtained from $H$ by deleting all hyperedges of size smaller than $m$, and denote by $H_{=m}$ the hypergraph consisting of the hyperedges of $H$ with size exactly $m$. Similarly, for a hypergraph family $\mathcal{H}$ let $\HH_m=\HH_{\ge m}=\{H_{\ge m}:H\in\HH\}$ and $\HH_{=m}=\{H_{=m}:H\in\HH\}$.
\end{definition}

\begin{definition} Let $\mathcal{H}$ be a hypergraph family. Denote by $m_\mathcal{H}(k)$ the smallest positive integer $m$ such that for every $H\in\mathcal{H}$ there exists a polychromatic $k$-coloring of the hypergraph $H_{\ge m}$. If there is no such $m$, set $m_{\mathcal{H}}(k)=\infty$.
\end{definition}

\begin{definition} Denote by $\A$ the hypergraph family which consists of those hypergraphs $H=(V,E)$ for which $V\subset \R^2$ is finite, and the edges are induced by axis-parallel strips: every $e\in E$ edge can be presented in the form $V\cap\{(x,y)\in\R^2: x_0<x<x_1\}$ or $V\cap\{(x,y)\in\R^2: y_0<y<y_1\}$ for some $x_0, x_1, y_0, y_1 \in\R$. Denote by $\A^*$ its dual: the family of such hypergraphs that the vertices are axis-parallel strips, and the edges are sets consisting of the strips containing a given point $(x,y)$. 
\end{definition}

\begin{definition} Denote by $\A_s$ the hypergraph family which consists of such $H=(V,E)$ hypergraphs for which $V\subset \R^2$ is finite, and the edges are induced by the union of $s$ axis-parallel strips: every $e\in E$ edge can be written in the form $V\cap \Big(\cup_{i=1}^s A_i\Big)$ where $A_1,\ldots, A_s$ are axis-parallel strips. 
\end{definition}

\begin{theorem}\label{s_strip}
$$m_{\A_s}(k)=s \cdot m_\A(k)-s+1.$$
\end{theorem}

\begin{proof}[Proof of Theorem~\ref{s_strip}]

For given $k\ge 2$ take a hypergraph from $\A$ which has a non-poly\-chromatic edge of size at least $m_\A(k)-1$ for any coloring. Take $k(s-1)+1$ copies of the point set that forms the vertices of this hypergraph, and place them along a diagonal line so that all edges of the original hypergraph would be still an edge. For any coloring of the resulting hypergraph, all copies have an edge of size $m_\A(k)-1$ that is not polychromatic, so there is a color which does not appear on the vertices of this hyperedge. By the pigeonhole principle we have $s$ of these hyperedges with the same missing color, these points are contained in the union of $s$ axis-parallel strips, so $s \cdot m_\A(k) -s + 1 \le m_{\A_s}(k)$.

For the other inequality we have to show the existence of a coloring such that any edge of size at least $s \cdot m_\A(k) -s + 1$ contains all $k$ colors. The same coloring as for $m_\A(k)$ satisfies this, because any set of $s \cdot m_\A(k) -s + 1$ points which is contained in the union of $s$ strips has a subset of $m_\A(k)$ points contained in a single strip.
\end{proof}

\end{document}
